\documentclass[a4paper,10pt]{article}
\usepackage[utf8]{inputenc}
\usepackage[ngerman]{babel}
\usepackage[margin=1.8cm]{geometry}
\usepackage{helvet}
\renewcommand{\familydefault}{\sfdefault}
\usepackage{tabularx}
\usepackage{booktabs}
\usepackage{amssymb}
\usepackage{wasysym}
\usepackage{xcolor}
\usepackage{tcolorbox}
\usepackage{enumitem}

\definecolor{primary}{RGB}{44, 62, 80}
\definecolor{accent}{RGB}{76, 175, 80}
\definecolor{lightgray}{RGB}{245, 247, 248}

\pagestyle{empty}

\newcommand{\radio}{\textsf{\Large$\bigcirc$}}
\newcommand{\checkbox}{\textsf{\Large$\square$}}
\newcommand{\linefill}{\rule{\linewidth}{0.4pt}\vspace{1.2ex}}

\begin{document}

% Header
{\Large \bfseries \color{primary} Evaluationsfragebogen – NWT Lernmodule}\\[2pt]
{\large \bfseries \color{accent} Fotometer \& Wetterstation-App | Piloterprobung}\\[8pt]

\begin{tcolorbox}[colback=lightgray,colframe=accent,arc=2mm,boxrule=1pt,top=6pt,bottom=6pt,left=8pt,right=8pt]
\textbf{Hinweis für Teilnehmende:} Dieser Fragebogen ist anonym. Deine ehrliche Rückmeldung hilft dabei, die Lernmodule zu verbessern. Es gibt keine richtigen oder falschen Antworten.\\[2pt]
\textbf{Zeitaufwand:} ca. 15–20 Minuten
\end{tcolorbox}

\vspace{4pt}

% Teil A
{\large \bfseries \color{primary} Teil A – Allgemeines \& Vorwissen}\\[-4pt]
\color{accent}\rule{\linewidth}{1.5pt}\color{black}\vspace{4pt}

\textbf{A1.} Welches Modul hast du heute bearbeitet? \textit{(Mehrfachauswahl möglich)}\\[4pt]
\checkbox~Fotometer (LDR \& Sensorkalibrierung) \quad 
\checkbox~Wetterstation (ESP32 \& IoT) \quad 
\checkbox~Beide

\vspace{8pt}

\textbf{A2.} Wie schätzt du dein Vorwissen in folgenden Bereichen \textbf{vor} der heutigen Sitzung ein?\\[4pt]
\begin{tabularx}{\linewidth}{X c c c c}
\toprule
\textbf{Bereich} & \textbf{Kein Vorwissen} & \textbf{Wenig} & \textbf{Einiges} & \textbf{Viel} \\
\midrule
Elektronik / Schaltkreise & \radio & \radio & \radio & \radio \\
Arduino / Microcontroller programmieren & \radio & \radio & \radio & \radio \\
Sensoren \& Messdaten & \radio & \radio & \radio & \radio \\
WLAN / Internet of Things & \radio & \radio & \radio & \radio \\
Datenauswertung (z. B. Excel, Tabellen) & \radio & \radio & \radio & \radio \\
\bottomrule
\end{tabularx}

\vspace{10pt}

% Teil B
{\large \bfseries \color{primary} Teil B – Verständlichkeit und Inhalt}\\[-4pt]
\color{accent}\rule{\linewidth}{1.5pt}\color{black}\vspace{2pt}
{\small\itshape Skala: 1 = trifft gar nicht zu $\cdot$ 5 = trifft voll zu}\\[4pt]

\textbf{B1.} Die Lernziele des Moduls waren für mich klar erkennbar.\\
\radio~1 \quad \radio~2 \quad \radio~3 \quad \radio~4 \quad \radio~5

\vspace{5pt}
\textbf{B2.} Die Erklärungen und Texte waren verständlich geschrieben.\\
\radio~1 \quad \radio~2 \quad \radio~3 \quad \radio~4 \quad \radio~5

\vspace{5pt}
\textbf{B3.} Die Fachbegriffe (z. B. LDR, Spannungsteiler, DHT11, IoT) wurden ausreichend erklärt.\\
\radio~1 \quad \radio~2 \quad \radio~3 \quad \radio~4 \quad \radio~5

\vspace{5pt}
\textbf{B4.} Der Schwierigkeitsgrad des Moduls war angemessen für mich.\\
\radio~1 (viel zu leicht) \quad \radio~2 \quad \radio~3 (genau richtig) \quad \radio~4 \quad \radio~5 (viel zu schwer)

\vspace{5pt}
\textbf{B5.} Die Schrittfolge (Abschnitte / Aufgaben) war logisch und nachvollziehbar aufgebaut.\\
\radio~1 \quad \radio~2 \quad \radio~3 \quad \radio~4 \quad \radio~5

\vspace{5pt}
\textbf{B6.} \textit{(Nur für Fotometer-Teilnehmende)} Das Konzept der Kalibrierung (Rohwert $\rightarrow$ Prozentwert) habe ich nach diesem Modul verstanden.\\
\radio~1 \quad \radio~2 \quad \radio~3 \quad \radio~4 \quad \radio~5 \quad \radio~Nicht zutreffend

\vspace{5pt}
\textbf{B7.} \textit{(Nur für Wetterstation-Teilnehmende)} Den Datenfluss vom Sensor über den ESP32 bis zur App (EVA-Prinzip / IoT-Kette) habe ich nach diesem Modul verstanden.\\
\radio~1 \quad \radio~2 \quad \radio~3 \quad \radio~4 \quad \radio~5 \quad \radio~Nicht zutreffend

\vspace{5pt}
\textbf{B8.} Welcher Teil des Moduls war für dich \textbf{am schwierigsten zu verstehen?} \textit{(offene Antwort)}\\
\linefill
\linefill

\newpage

% Teil C
{\large \bfseries \color{primary} Teil C – Digitale Elemente \& Interaktivität}\\[-4pt]
\color{accent}\rule{\linewidth}{1.5pt}\color{black}\vspace{2pt}
{\small\itshape Skala: 1 = überhaupt nicht hilfreich $\cdot$ 5 = sehr hilfreich}\\[4pt]

\textbf{C1.} Wie hilfreich waren folgende Elemente für dein Verständnis?\\[4pt]
\begin{tabularx}{\linewidth}{X c c c c c c}
\toprule
\textbf{Element} & \textbf{1} & \textbf{2} & \textbf{3} & \textbf{4} & \textbf{5} & \textbf{Nicht genutzt} \\
\midrule
Interaktive Simulationen / Schieberegler & \radio & \radio & \radio & \radio & \radio & \radio \\
Code-Beispiele mit Kopier-Button & \radio & \radio & \radio & \radio & \radio & \radio \\
Tooltip-Erklärungen (Hover-Infos) & \radio & \radio & \radio & \radio & \radio & \radio \\
Flip-Karten (Klick zum Umdrehen) & \radio & \radio & \radio & \radio & \radio & \radio \\
Diagramme / Kennlinien-Grafiken & \radio & \radio & \radio & \radio & \radio & \radio \\
Schaltplan / Fritzing-Abbildungen & \radio & \radio & \radio & \radio & \radio & \radio \\
Drag-and-Drop Reihenfolge-Übung (Modul 2) & \radio & \radio & \radio & \radio & \radio & \radio \\
Lernstandsquiz / Wissensfragen & \radio & \radio & \radio & \radio & \radio & \radio \\
Fortschrittsbalken oben & \radio & \radio & \radio & \radio & \radio & \radio \\
\bottomrule
\end{tabularx}

\vspace{8pt}

\textbf{C2.} Welches dieser Elemente hat dir am \textbf{meisten} geholfen? Warum?\\
\linefill
\linefill

\vspace{4pt}
\textbf{C3.} Gab es digitale Elemente, die dich \textbf{verwirrt oder gestört} haben?\\
\linefill
\linefill

\vspace{8pt}

% Teil D
{\large \bfseries \color{primary} Teil D – Praktische Durchführung}\\[-4pt]
\color{accent}\rule{\linewidth}{1.5pt}\color{black}\vspace{4pt}

\textbf{D1.} Der Aufbau der Hardware (Breadboard, Verkabelung) hat gut funktioniert.\\
\radio~1 (gar nicht) \quad \radio~2 \quad \radio~3 \quad \radio~4 \quad \radio~5 (sehr gut) \quad \radio~Hatte technische Probleme

\vspace{5pt}
\textbf{D2.} Die Anleitungen im Modul haben mir geholfen, die Hardware korrekt aufzubauen.\\
\radio~1 \quad \radio~2 \quad \radio~3 \quad \radio~4 \quad \radio~5

\vspace{5pt}
\textbf{D3.} Das Hochladen des Codes auf den Microcontroller hat geklappt.\\
\radio~Ja, ohne Probleme \quad \radio~Ja, mit Hilfe \quad \radio~Nein, trotz Hilfe nicht \quad \radio~Nicht versucht

\vspace{5pt}
\textbf{D4.} Was wäre hilfreich gewesen, um technische Probleme schneller zu lösen?\\
\linefill
\linefill

\vspace{5pt}
\textbf{D5.} Wie lange hast du (ungefähr) für das Modul gebraucht?\\
\radio~$< 30$ Min \quad \radio~$30$--$60$ Min \quad \radio~$60$--$90$ Min \quad \radio~$> 90$ Min

\vspace{8pt}

% Teil E
{\large \bfseries \color{primary} Teil E – Lernertrag \& Motivation}\\[-4pt]
\color{accent}\rule{\linewidth}{1.5pt}\color{black}\vspace{4pt}

\textbf{E1.} Ich habe heute etwas gelernt, das ich vorher nicht wusste oder konnte.\\
\radio~1 (gar nicht) \quad \radio~2 \quad \radio~3 \quad \radio~4 \quad \radio~5 (sehr viel gelernt)

\vspace{5pt}
\textbf{E2.} Das Thema des Moduls ist für mich relevant und interessant.\\
\radio~1 \quad \radio~2 \quad \radio~3 \quad \radio~4 \quad \radio~5

\vspace{5pt}
\textbf{E3.} Ich sehe eine hohe Relevanz und Wichtigkeit dieser Inhalte für meinen zukünftigen Lehrerberuf (z. B. für den späteren NwT-Unterricht).\\
\radio~1 \quad \radio~2 \quad \radio~3 \quad \radio~4 \quad \radio~5

\vspace{5pt}
\textbf{E4.} Das Modul hat mein Interesse für Elektronik / Programmierung / IoT geweckt oder gesteigert.\\
\radio~1 \quad \radio~2 \quad \radio~3 \quad \radio~4 \quad \radio~5

\vspace{5pt}
\textbf{E5.} Ich könnte die Grundprinzipien aus diesem Modul jemandem erklären, der es nicht kennt.\\
\radio~Ja, ich bin sicher \quad \radio~Wahrscheinlich schon \quad \radio~Nur teilweise \quad \radio~Eher nicht

\vspace{5pt}
\textbf{E6.} Was nimmst du als \textbf{wichtigste Erkenntnis} aus dem heutigen Tag mit?\\
\linefill
\linefill
\linefill

\vspace{8pt}

% Teil F
{\large \bfseries \color{primary} Teil F – Gestaltung \& Nutzererfahrung}\\[-4pt]
\color{accent}\rule{\linewidth}{1.5pt}\color{black}\vspace{4pt}

\textbf{F1.} Das Layout und Design der Webseite war übersichtlich und angenehm zu benutzen.\\
\radio~1 \quad \radio~2 \quad \radio~3 \quad \radio~4 \quad \radio~5

\vspace{5pt}
\textbf{F2.} Die Webseite hat auf meinem Gerät / Browser einwandfrei funktioniert.\\
\radio~Ja \quad \radio~Teilweise (mit Problemen) \quad \radio~Nein \quad \textit{Gerät/Browser:} \rule{4cm}{0.4pt}

\vspace{5pt}
\textbf{F3.} Was würdest du am Design oder der Benutzerführung \textbf{verbessern}?\\
\linefill
\linefill

\vspace{8pt}

% Teil G
{\large \bfseries \color{primary} Teil G – Gesamtbewertung \& offene Rückmeldung}\\[-4pt]
\color{accent}\rule{\linewidth}{1.5pt}\color{black}\vspace{4pt}

\textbf{G1.} Gesamtbewertung des Lernmoduls:\\
\radio~1 (sehr schlecht) \quad \radio~2 \quad \radio~3 \quad \radio~4 \quad \radio~5 (ausgezeichnet)

\vspace{5pt}
\textbf{G2.} Würdest du dieses Modul einem Kommilitonen / einer Kommilitonin weiterempfehlen?\\
\radio~Ja, auf jeden Fall \quad \radio~Ja, eher schon \quad \radio~Unentschieden \quad \radio~Eher nicht \quad \radio~Nein

\vspace{5pt}
\textbf{G3.} Was hat dich am Modul am \textbf{meisten überzeugt}?\\
\linefill
\linefill

\vspace{5pt}
\textbf{G4.} Was sollte deiner Meinung nach \textbf{unbedingt verbessert oder ergänzt} werden?\\
\linefill
\linefill
\linefill

\vspace{5pt}
\textbf{G5.} Hast du weitere Anmerkungen, Ideen oder Wünsche?\\
\linefill
\linefill

\vspace{10pt}
\centering{\textbf{Vielen Dank für deine Zeit und dein Feedback!}}

\vspace{12pt}
\begin{tcolorbox}[colback=lightgray,colframe=gray,arc=1mm,boxrule=0.5pt,top=4pt,bottom=4pt,left=6pt,right=6pt]
\small \textbf{Für den Auswertenden (intern):}\\
$\bullet$ Teil B: Verständlichkeit $\rightarrow$ Mittelwert bilden, Freitexte qualitativ kodieren\\
$\bullet$ Teil C: Matrix $\rightarrow$ Identifikation besonders wirksamer / störender Elemente\\
$\bullet$ Teil D: Hardware-Hürden $\rightarrow$ Direkt als Verbesserungsmaßnahmen verwertbar\\
$\bullet$ Teil E: Lernertrag $\rightarrow$ B4 + E1 + E3 + E5 kombiniert für Lernertrag-Relevanz-Matrix\\
$\bullet$ Teil G: Freitexte $\rightarrow$ Inhaltsanalyse nach Kategorien (Inhalt / Design / Technik / Motivation)
\end{tcolorbox}

\end{document}
